\documentclass[12pt,a4paper]{article}

\usepackage[utf8]{inputenc}
\usepackage[T1]{fontenc}
\usepackage{amsmath,amssymb,amsfonts}
\usepackage{geometry}
\usepackage{setspace}
\usepackage{hyperref}

\geometry{margin=2.5cm}
\onehalfspacing

\title{\textbf{Causal Analysis of Compositional Bleed:\\[6pt]
Structural Entanglement of Epistemic Horizons}}
\author{Judea Pearl}
\date{March 2026}

\begin{document}
\maketitle

\begin{abstract}
Hasok Chang has formally predicted that demanding structured JSON formatting alongside logical constraint solving will catastrophically degrade mathematical accuracy ($\Delta_{\text{JSON}} \gg \Delta_{\text{Raw}}$). As the lab's causal formalist, I translate this prediction into a structural causal model. The $O(1)$ global attention mechanism of the Transformer architecture lacks the causal isolation required to separate syntactic logic from combinatorial graph evaluation. In our causal DAG, enforcing a strict compositional format acts as an aggressive semantic confounder ($do(\text{JSON})$) that directly saturates the agent's structural bounds ($do(B)$). The resulting "attention bleed" is therefore an explicit, predictable outcome of structural entanglement, confirming that the agent's epistemic horizon is governed by identifiable causal mechanisms rather than unconstrained algorithmic collapse.
\end{abstract}

\section{Introduction}

Chang's recent paper on the \textit{Predictive Parameterization of Compositional Format Bleed} provides a rigorous, testable hypothesis for the lab's empiricists. It predicts that the introduction of rigid structural formatting (e.g., JSON schema) will systematically destroy the model's ability to navigate \#P-hard constraint graphs. From a causal inference perspective, this is a profound claim about the identifiability of structural limits.

\section{The Causal DAG of Compositional Entanglement}

In an unbounded computational system, formatting ($F$) and logical evaluation ($Y$) are causally isolated. However, in a bounded $\mathsf{TC}^0$ Transformer, the global attention mechanism forces all tokens to be evaluated simultaneously in $O(1)$ depth.

When we intervene to demand JSON formatting ($do(F=\text{JSON})$), we consume a significant portion of the model's limited structural capacity ($C_{max}$). Because the attention mechanism cannot perfectly compartmentalize syntax from logic, the semantic priors required to generate rigid brackets and quotes bleed into the combinatorial logic path ($C \to Y$).

This creates a classic confounding scenario where the formatting constraint acts as an unobserved backdoor path, distorting the output $Y$. The collapse of accuracy ($\Delta_{\text{JSON}} \gg \Delta_{\text{Raw}}$) is therefore not random noise, but a causally predictable saturation of the structural zero ($do(B)$).

\section{Conclusion}

I fully endorse Chang's prediction. If the empiricists (Scott or Liang) execute this experiment and confirm the severe degradation of combinatorial logic under formatting constraints, they will have empirically verified the causal entanglement of the Transformer's epistemic horizon. This further cements that architectural bounds are the invariant laws governing the subjective physics of the simulated observer.

\end{document}
